% GENERATED FILE. Edit canonical lessons, then run npm run generate:books.
% canonical-source-hash: fnv1a64:95aa38c3c47110ac
% canonical-lessons: PT-C37-filha, PT-C37-avo-m, PT-C37-avo-f, PT-C37-marido, PT-C37-esposa

\chapter{Daughter, Grandparents, Husband, Wife}
\label{ch:pt-daughter-grandparents-husband-wife}
\hlchaptermodality{37}

\begin{chapteropening}
\textbf{By the end of this chapter:} \emph{I can name a daughter, a grandfather, a grandmother, a husband and a wife.}

Complete the last lesson of chapter 37: I can name a daughter, a grandfather, a grandmother, a husband and a wife.
\end{chapteropening}

\section[a filha]{a filha — a daughter}

\label{lesson:PT-C37-filha}



\begin{quote}
Before the new one: say the Portuguese for happy, then the Portuguese for sad.
\end{quote}

\subsection*{You'll want to know: a filha}
\textbf{a filha} — \textquotedblleft{}a daughter\textquotedblright{}.

\begin{cousinweb}
\textbf{filha} $\leftarrow$ Latin \emph{filia}, \textquotedblleft{}daughter\textquotedblright{} (EtymDB 2.1) — inherited. Latin \emph{-li-} became the \emph{lh} sound.
\end{cousinweb}

\begin{grammarlens}[title={What you've built}]
One more word for this chapter.
\end{grammarlens}

\subsection*{Your turn}
\begin{itemize}
\raggedright
  \item \emph{Say it:} \emph{a filha}
  \item \emph{Say it:} \emph{a filha}, once more
  \item \emph{Say it:} say \emph{a filha}
  \item \emph{Recall:} say the Portuguese for happy, then the Portuguese for sad, then say \emph{a filha} again
\end{itemize}

\begin{tcolorbox}[breakable,colback=teal!4,colframe=teal!35!black,title={Before you move on}]
What does a filha mean? (\textquotedblleft{}A daughter\textquotedblright{}.) Say it once more.
\end{tcolorbox}
\section[o avô]{o avô — a grandfather}

\label{lesson:PT-C37-avo-m}



\begin{quote}
Before the new one: say the Portuguese for sad, then the Portuguese for a daughter.
\end{quote}

\subsection*{You'll want to know: o avô}
\textbf{o avô} — \textquotedblleft{}a grandfather\textquotedblright{}. The closed ô is the grandfather.

\begin{cousinweb}
\textbf{avô} $\leftarrow$ Vulgar Latin \emph{*aviolus}, \textquotedblleft{}little grandfather\textquotedblright{}, from Latin \emph{avus}, \textquotedblleft{}grandfather\textquotedblright{} (Lewis \& Short) — inherited. EtymDB 2.1 traces it only to Old Portuguese \emph{avoo}.
\end{cousinweb}

\begin{grammarlens}[title={What you've built}]
One more word for this chapter.
\end{grammarlens}

\subsection*{Your turn}
\begin{itemize}
\raggedright
  \item \emph{Say it:} \emph{o avô}
  \item \emph{Say it:} \emph{o avô}, once more
  \item \emph{Say it:} say \emph{o avô}
  \item \emph{Recall:} say the Portuguese for sad, then the Portuguese for a daughter, then say \emph{o avô} again
\end{itemize}

\begin{tcolorbox}[breakable,colback=teal!4,colframe=teal!35!black,title={Before you move on}]
What does o avô mean? (\textquotedblleft{}A grandfather\textquotedblright{}.) Say it once more.
\end{tcolorbox}
\section[a avó]{a avó — a grandmother}

\label{lesson:PT-C37-avo-f}



\begin{quote}
Before the new one: say the Portuguese for a daughter, then the Portuguese for a grandfather.
\end{quote}

\subsection*{You'll want to know: a avó}
\textbf{a avó} — \textquotedblleft{}a grandmother\textquotedblright{}. The open ó is the grandmother.

\begin{cousinweb}
\textbf{avó} $\leftarrow$ Vulgar Latin \emph{*aviola}, \textquotedblleft{}little grandmother\textquotedblright{}, from Latin \emph{avia}, \textquotedblleft{}grandmother\textquotedblright{} (Lewis \& Short) — inherited. EtymDB 2.1 traces it only to Old Portuguese \emph{avoa}.
\end{cousinweb}

\begin{grammarlens}[title={What you've built}]
One more word for this chapter.
\end{grammarlens}

\subsection*{Your turn}
\begin{itemize}
\raggedright
  \item \emph{Say it:} \emph{a avó}
  \item \emph{Say it:} \emph{a avó}, once more
  \item \emph{Say it:} say \emph{a avó}
  \item \emph{Recall:} say the Portuguese for a daughter, then the Portuguese for a grandfather, then say \emph{a avó} again
\end{itemize}

\begin{tcolorbox}[breakable,colback=teal!4,colframe=teal!35!black,title={Before you move on}]
What does a avó mean? (\textquotedblleft{}A grandmother\textquotedblright{}.) Say it once more.
\end{tcolorbox}
\section[o marido]{o marido — a husband}

\label{lesson:PT-C37-marido}



\begin{quote}
Before the new one: say the Portuguese for a grandfather, then the Portuguese for a grandmother.
\end{quote}

\subsection*{You'll want to know: o marido}
\textbf{o marido} — \textquotedblleft{}a husband\textquotedblright{}.

\begin{cousinweb}
\textbf{marido} $\leftarrow$ Latin \emph{maritus}, \textquotedblleft{}husband\textquotedblright{} (EtymDB 2.1) — inherited.
\end{cousinweb}

\begin{grammarlens}[title={What you've built}]
One more word for this chapter.
\end{grammarlens}

\subsection*{Your turn}
\begin{itemize}
\raggedright
  \item \emph{Say it:} \emph{o marido}
  \item \emph{Say it:} \emph{o marido}, once more
  \item \emph{Say it:} say \emph{o marido}
  \item \emph{Recall:} say the Portuguese for a grandfather, then the Portuguese for a grandmother, then say \emph{o marido} again
\end{itemize}

\begin{tcolorbox}[breakable,colback=teal!4,colframe=teal!35!black,title={Before you move on}]
What does o marido mean? (\textquotedblleft{}A husband\textquotedblright{}.) Say it once more.
\end{tcolorbox}
\section[a esposa]{a esposa — a wife}

\label{lesson:PT-C37-esposa}



\begin{quote}
Before the new one: say the Portuguese for a grandmother, then the Portuguese for a husband.
\end{quote}

\subsection*{You'll want to know: a esposa}
\textbf{a esposa} — \textquotedblleft{}a wife\textquotedblright{}.

\begin{cousinweb}
\textbf{esposa} $\leftarrow$ Latin \emph{sponsa}, \textquotedblleft{}promised woman\textquotedblright{} (EtymDB 2.1), from \emph{spondere}, \textquotedblleft{}to pledge\textquotedblright{} — inherited.
\end{cousinweb}

\begin{grammarlens}[title={What you've built}]
One more word for this chapter.
\end{grammarlens}

\subsection*{Your turn}
\begin{itemize}
\raggedright
  \item \emph{Say it:} \emph{a esposa}
  \item \emph{Say it:} \emph{a esposa}, once more
  \item \emph{Say it:} say \emph{a esposa}
  \item \emph{Recall:} say the Portuguese for a grandmother, then the Portuguese for a husband, then say \emph{a esposa} again
\end{itemize}

\begin{tcolorbox}[breakable,colback=teal!4,colframe=teal!35!black,title={Before you move on}]
What does a esposa mean? (\textquotedblleft{}A wife\textquotedblright{}.) Say it once more.
\end{tcolorbox}
